Scale the circle radius to one and write the common angle $\beta$. The angle at $C^{\prime}$ subtended by the unit circle is at most $30^\circ$, since $OC^{\prime}=2$. Hence $0<\beta\le30^\circ$. Chord $AC$ has length $2\sin\beta$, and $\angle ACO=90^\circ-\beta$. Since $CC^{\prime}$ is opposite $CO$, $\angle ACC^{\prime}=90^\circ+\beta$, leaving $\angle CAC^{\prime}=90^\circ-2\beta$. The sine law in $ACC^{\prime}$ gives $\dfrac{CC^{\prime}}{\cos2\beta}=\dfrac{AC}{\sin\beta}=2$. Since $CC^{\prime}=1$, $\cos2\beta=\dfrac{1}{2}$, so $\boxed{\beta=30^\circ}$.
\begin{asy}
import graph;import geometry;usepackage("amsmath");size(7cm);
pair O=(0,0),C=(1,0),Cp=(2,0),A=dir(60),B=(-1,0);draw(circle(O,1));draw(A--B--C--A);draw(O--Cp--A);label("$O$",O,(0,-1));label("$C$",C,(0,-1));label("$C^{\prime}$",Cp,(1,0));label("$A$",A,(0,1));label("$B$",B,(-1,0));
\end{asy}

